% Generated draft: requires mathematical and source review.
\documentclass[11pt]{article}
\usepackage[margin=25mm]{geometry}
\usepackage{fontspec}
\setmainfont{Latin Modern Roman}
\usepackage{amsmath,amssymb,mathtools,unicode-math}
\setmathfont{Latin Modern Math}
\usepackage{enumitem,needspace,xurl,hyperref}
\setlength{\parindent}{0pt}
\setlength{\parskip}{5pt}
\setlength{\emergencystretch}{3em}
\newcommand{\N}{\mathbb{N}}
\newcommand{\Z}{\mathbb{Z}}
\newcommand{\Q}{\mathbb{Q}}
\newcommand{\R}{\mathbb{R}}
\newcommand{\C}{\mathbb{C}}
\newcommand{\F}{\mathbb{F}}
\newcommand{\A}{\mathbb{A}}
\newcommand{\PP}{\mathbb{P}}
\newcommand{\E}{\mathbb{E}}
\newcommand{\eps}{\varepsilon}
\newcommand{\dd}{\,\mathrm{d}}
\begin{document}
\section*{UnsolvedMath ID: 30001062}
\textbf{Source problem code:} catalogue entry
\subsection{Definitions and mathematical statement}
Let $q$ be a prime power, let $d\geq2$, and let $E\subseteq\F_q^d$. Let
\[
V=\{0,v_1,\ldots,v_{k-1}\}\subseteq\F_q^d
\]
be the source's nondegenerate $k$-point configuration, with $v_1,\ldots,v_{k-1}$ linearly independent. In particular $k-1\leq d$. The requested copy is specifically
\[
\tau+O(V)=\{\tau,\tau+Ov_1,\ldots,\tau+Ov_{k-1}\}\subseteq E,
\qquad \tau\in\F_q^d,\quad O\in\operatorname{SO}_d(\F_q).
\]
The special orthogonal group, not the full orthogonal group or a similarity group, is retained.

For a fixed admissible configuration, an exact threshold formulation is
\[
m(q,d,V)=\min\{m:\text{every }E\subseteq\F_q^d\text{ with }|E|\geq m
\text{ contains some }\tau+O(V)\}.
\]
Determine this threshold, or its intended sharp order as $q$ grows. If a bound depending only on $(q,d,k)$ is sought, it must hold uniformly over all configurations in the specified nondegenerate class. The record does not specify whether constants may depend on $V$ or the exact precision intended in ``how large.''

The quadratic form behind $\operatorname{SO}_d$ and the characteristic hypotheses must be fixed. In the usual odd-characteristic standard-form convention,
\[
\operatorname{SO}_d(\F_q)=\{O\in\operatorname{GL}_d(\F_q):O^{\mathsf T}O=I,\ \det O=1\}.
\]
This is explanatory, not an added odd-characteristic hypothesis. The record does not exclude characteristic two. Likewise, linear independence of the $v_i$ does not by itself state that their Gram matrix for the intended bilinear form is nonsingular. Whether ``nondegenerate'' includes that additional quadratic restriction requires source confirmation. No such restriction is guessed.

The supplied literature assessment checked 2026-08-21 labels the problem partially solved. It attributes substantial density thresholds to Hart--Iosevich and later work, for congruent copies or positive proportions of finite-field simplices, while leaving the optimal threshold for the full Problem 5 unresolved. Guaranteeing a positive proportion of congruence classes is not silently identified with containing every prescribed admissible configuration.
\subsection{Short English statement}
Determine the size threshold forcing a translated special-orthogonal copy of a prescribed independent finite-field configuration. Preserve the distinction between fixed-configuration and uniform thresholds, and between affine independence and quadratic nondegeneracy.
\subsection{Sources}
\begin{itemize}
\item Local UnsolvedMath registry, permanent ID 30001062, source code catalogue entry; supplied complete record, whose corrected-verified note restores the finite-field configuration notation. The assessment refers to the report's Problem 5. Literature checked 2026-08-21: the assessment is supplied dated curation, not independently verified here.
\item Discrete Geometry (2009), Oberwolfach Reports, DOI 10.4171/owr/2008/44, pp. 2477--2556: \url{https://doi.org/10.4171/owr/2008/44}. Proposed-year metadata gives 2008.
\item Chapman--Erdogan--Hart--Iosevich--Koh: \url{https://arxiv.org/abs/0903.4218}. The supplied progress assessment also names Hart--Iosevich without a separate bibliographic entry.
\end{itemize}
\end{document}
